% Pacotes utilizados
\usepackage[utf8]{inputenc} % Codificação dos caracteres
\usepackage[T1]{fontenc}    % Melhor suporte a acentuação
\usepackage{hyperref}       % Permite criar links
\usepackage{graphicx}       % Para incluir imagens
\usepackage{amssymb}        % Símbolos matemáticos
\usepackage{amsmath}        % Expressões matemáticas
\usepackage{ragged2e}       % Justificar textos
\usepackage{array}          % Matrizes e tabelas avançadas
%\usepackage[shorthands=off,portuguese]{babel}  % Idioma português do Brasil
\usepackage{verbatim}       % Bloco de comentários
\usepackage{booktabs}       % Tabelas elegantes
\usepackage{tikz}           % Desenhos e diagramas TikZ
\usetikzlibrary{positioning, arrows.meta}
\usepackage[ruled, vlined, linesnumbered, slide]{algorithm2e} % Pseudocódigo para slides

